\section{Koinzidenzeinheit}

Eine Koinzidenzeinheit wird genutzt um das gleichzeitige Auftreten von Signalen zu überprüfen. Als Eingang werden üblicherweise Logiksignale genutzt und im Falle von Koinzidenz wird ein Logiksignal erzeugt. Diese Funktionalität ist äquivalent zu einem logischen Und\cite[295]{leotechniques}.

Die Verwendete Koinzidenzeinheit ist das Modell 418A der Firma \texttt{ORTEC}, welche mit bis zu vier Signalen mit Mindestpulsdauern von $\SI{50}{\nano\second}$ arbeiten kann und im Falle einer Koinzidenz Logiksignale mit einer Pulsdauer von $\SI{500}{\nano\second}$ erzeugt. Die Auflösungszeit des Kanals A wird hierbei über ein Potentiometer von $\SIrange{0.1}{2}{\micro\second}$ eingestellt und bei allen anderen Kanälen durch die Pulsdauer festgelegt\cite{coincidence_unit}.
